\section*{Chapter \ref{ch:b2:affine} --- Affine Spaces}

\begin{solution}{exo:b2:affine:1}
$\{x + y + z = 1\}$: affine plane, direction the vector plane $\{x
+ y + z = 0\}$, dimension $2$. Adding $x - z = 3$: an affine line
(two independent equations), direction $\{x + y + z = 0,\ x = z\}
= \operatorname{Vect}\bigl((1, -2, 1)\bigr)$, dimension $1$. $\{x^2
+ y^2 = 1\}$: a cylinder --- not stable under barycenters (the
midpoint of $(1,0,0)$ and $(-1,0,0)$ is the origin, off the
cylinder): not affine. A compatible system $MX = B$: affine
subspace $X_0 + \ker M$ of dimension $\dim\ker M$, as recalled in
\cref{def:b2:affine:subspace}.
\end{solution}

\begin{solution}{exo:b2:affine:2}
$C = \operatorname{bar}(A, 1-t; B, t)$ means $\vect{AC} =
t\,\vect{AB}$: existence of such $t$ is exactly linkage of
$\vect{AC}$ with $\vect{AB} \neq 0$, i.e.\ alignment. Positions: $t
= \frac12$: midpoint; $t = 2$: beyond $B$, at $B$'s distance from
it ($\vect{AC} = 2\vect{AB}$); $t = -1$: the reflection of $B$
through $A$.
\end{solution}

\begin{solution}{exo:b2:affine:3}
$M$ is the projection matrix onto $\operatorname{Vect}(1,1)$ along
$(1,-1)$ (check $M^2 = M$). Image of $f$: $\{MX + C\} = C +
\operatorname{im} M$: the affine line through $(1,0)$ directed by
$(1,1)$. Fixed points: $X = MX + C$, i.e.\ $(I - M)X = C$; but $C =
(1, 0)^{\mathsf T}$ and $\operatorname{im}(I - M) =
\operatorname{Vect}(1,-1)$; is $(1,0)$ in it? $(1, 0) =
\alpha(1,-1)$ forces $\alpha = 1$ and $0 = -1$: no. No fixed
points. And
\[
f(f(X)) = M(MX + C) + C = MX + MC + C = f(X) + MC,
\qquad MC = \tfrac12(1,1)^{\mathsf T} :
\]
$f\circ f$ is $f$ followed by a translation along the image line
--- $f$ is a ``glide projection'': projection onto the line
composed with a slide.
\end{solution}

\begin{solution}{exo:b2:affine:4}
$I = \operatorname{bar}(B, 2; C, 1)$ (since $\vect{BI} =
\frac13\vect{BC}$ places $I$ closer to $B$: weights $2$ on $B$,
$1$ on $C$ --- check: $\vect{BI} = \frac{1}{3}\vect{BC}$).
Similarly $J = \operatorname{bar}(C, 2; A, 1)$, $K =
\operatorname{bar}(A, 2; B, 1)$. Summing the three weighted
systems, the barycenter of $(I, 1; J, 1; K, 1)$ (each of total
weight $3$, so replace $I$ by its system, etc.) is
\[
\operatorname{bar}\bigl(A, 1 + 2;\ B, 2 + 1;\ C, 1 + 2\bigr)
= \operatorname{bar}(A, 1; B, 1; C, 1) = G ,
\]
the centroid of $ABC$: the triangle $IJK$ has the same centroid ---
predictable, because the construction treats $A, B, C$ cyclically
and the centroid is the unique fixed point of the cyclic symmetry
of weights.
\end{solution}

\begin{solution}{exo:b2:affine:5}
Set $g(u) = f(O + u) - f(O)$ (working in $\R^n$ vectorialized at
$O$), $g(0) = 0$.

\emph{Additivity:} $\frac{(O + u) + (O + v)}{2} = O + \frac{u +
v}{2}$, so midpoint preservation gives $g\bigl(\frac{u+v}{2}\bigr)
= \frac{g(u) + g(v)}{2}$; with $v = 0$: $g(u/2) = g(u)/2$;
combining, $g(u + v) = 2g\bigl(\frac{u+v}{2}\bigr) = g(u) + g(v)$.

\emph{$\Q$-homogeneity:} additivity gives $g(nu) = ng(u)$ ($n \in
\N$, induction), then $g(-u) = -g(u)$ (add), then $g(\frac pq u) =
\frac pq g(u)$ (apply $q$, use injectivity of scaling).

\emph{$\R$-homogeneity:} for $t \in \R$, take rationals $t_n \to
t$: $g(t_nu) = t_ng(u)$, and continuity of $g$ (inherited from $f$)
passes to the limit: $g(tu) = tg(u)$. Hence $g$ is linear and $f =
f(O) + g$: affine.
\end{solution}

\begin{solution}{exo:b2:affine:6}
Linear part $\varphi(x,y) = (y, x)$: the reflection in the diagonal
$y = x$ (orthogonal, determinant $-1$). Fixed points: $(x, y) = (y
+ 1, x - 1)$ amounts to the single equation $y = x - 1$ (the two
components are equivalent): every point of the line $y = x - 1$ is
fixed. So $f$ fixes that line pointwise: $f$
is the \emph{reflection} in that axis (an isometry with a line of
fixed points and linear part a reflection). Consistently, $f \circ
f(x,y) = f(y+1, x-1) = (x - 1 + 1, y + 1 - 1) = (x, y)$: an
involution, as a reflection must be.
\end{solution}

\begin{solution}{exo:b2:affine:7}
The $n + 1$ vectors $\vect{A_1A_i}$ ($i \geq 2$) are linked in
dimension $n$: there are $\mu_i$, not all zero, with $\sum_{i\geq2}
\mu_i\vect{A_1A_i} = 0$; set $\mu_1 = -\sum_{i \geq 2}\mu_i$, so
$\sum_{i}\mu_i = 0$ and $\sum_i \mu_i\,\vect{OA_i} = 0$ for every
$O$, with not all $\mu_i$ zero. Split indices: $P = \{i : \mu_i >
0\}$, $N = \{i : \mu_i < 0\}$, both nonempty (the $\mu_i$ sum to
zero and are not all zero). With $s = \sum_{i\in P}\mu_i =
-\sum_{i \in N}\mu_i > 0$:
\[
\operatorname{bar}\bigl(A_i, \tfrac{\mu_i}{s}\bigr)_{i \in P}
= \operatorname{bar}\bigl(A_i, \tfrac{-\mu_i}{s}\bigr)_{i \in N},
\]
(both sides equal the point $X$ with $\vect{OX} = \frac1s\sum_{i\in
P}\mu_i\vect{OA_i}$, by the relation): a common point of the two
convex hulls, with disjoint index groups.
\end{solution}

\begin{solution}{exo:b2:affine:8}
\emph{Image = fixed points:} for $Y = f(X)$, $f(Y) = f(f(X)) =
f(X) = Y$: every image point is fixed; conversely fixed points are
images. So $\mathcal{F} = \operatorname{im} f =
\operatorname{Fix}(f)$ is nonempty, and it is an affine subspace
(image of an affine map), with direction $\operatorname{im}\vec f$.

\emph{Projection structure:} $\vec f$ is idempotent ($\vec{f\circ
f} = \vec f^{\,2} = \vec f$), so $E = \operatorname{im}\vec f
\oplus \ker \vec f$ (\cref{ex:b2:reduction:projections}). For any
point $X$, consider the vector $\vect{f(X)\,X}$; applying $\vec f$:
\[
\vec f\bigl(\vect{f(X)\,X}\bigr) = \vect{f(f(X))\,f(X)} = 0
\qquad (f \circ f = f),
\]
so $\vect{f(X)\,X} \in \ker\vec f$. Hence $X = f(X) +
\vect{f(X)X}$ displays $X$ as a point of $\mathcal{F}$ translated
by a vector of $\ker\vec f$: $f$ is exactly the projection onto
$\mathcal{F}$ along $\ker\vec f$. Conversely such projections
clearly satisfy $f \circ f = f$.
\end{solution}
